\SourceLeaf{378}
\par\bigskip\Needspace{9\baselineskip}
\pdfbookmark[1]{Theorie der biquadratischen Reste. VI.}{nachlass-biquadratic-VI}
\centerline{[VI.]}\medskip
\begin{center}\large\scshape Theorie der biquadratischen Reste.\end{center}
\BookRule{.18\textwidth}
\renewcommand{\Article}[1]{\par\medskip\Needspace{4\baselineskip}\pdfbookmark[2]{Art. #1}{nachlass-biquadratic-VI-art-#1}\centerline{\textbf{#1.}}\nobreak\smallskip}
\Article{1}
Eine Reihe ganzer complexer Zahlen $\varphi,\varphi',\varphi''$ u.\,s.\,w. sei so beschaffen, dass erstlich sie unter einander alle incongruent sind nach dem Modulus $\mu=\alpha+\beta i$, $\alpha$ und $\beta$ ganze reelle Zahlen bezeichnend, zweitens dass jede ganze complexe Zahl einer von jenen nach dem Modulus $\mu$ congruent ist. Unter dieser Voraussetzung heisst der Inbegriff der Zahlen $\varphi,\varphi',\varphi''$ u.\,s.\,w. ein vollständiges Restsystem für den Modulus $\mu$. Es ist bewiesen, dass die Anzahl der darin begriffenen Zahlen der Norm von $\mu$, d.\,i. der Zahl $\alpha\alpha+\beta\beta$ gleich ist, welche mit $\nu$ bezeichnet werden soll.

\Article{2}
Unter den Zahlen $\varphi,\varphi',\varphi''$ u.\,s.\,w. ist Eine durch $\mu$ theilbare; wird dieselbe ausgeschlossen und der Inbegriff der übrigen mit $\chi$ bezeichnet, so bildet $\chi$ ein vollständiges System der durch den Modulus untheilbaren Reste, deren Anzahl $=\nu-1$. Beschränken wir die Untersuchung auf ungerade Modulen, so ist $\nu-1$ durch $4$ theilbar. Es werden dann ferner $f,if,-f,-if$ unter sich incongruent sein, folglich diejenigen Zahlen in $\chi$, welche resp.\ denen $if,-f,-if$ congruent sind, unter sich und von $f$ verschieden. (Associirte und zusammengesetzte Zahlen.)

Hieraus ergiebt sich eine Zerlegung von $\chi$ in vier Gruppen oder partielle Systeme $\kappa,\kappa',\kappa'',\kappa'''$. Man setzt eine beliebige Zahl aus $\chi$, z.\,B.\ $\varphi$ in die Gruppe $\kappa$, und die drei den Zahlen $i\varphi,-\varphi,-i\varphi$ congruenten Glieder von $\chi$, der Reihe nach in die Gruppen $\kappa',\kappa'',\kappa'''$. Nachdem diese vier Zahlen aus $\chi$ gestrichen sind, setzt man eine beliebige der übrigen wieder in $\kappa$, und die drei auf ähnliche Art davon abhängigen in $\kappa',\kappa'',\kappa'''$. So fährt man fort, bis das ganze System $\chi$ vertheilt ist. Die Gruppen $\kappa,\kappa',\kappa'',\kappa'''$ sollen zusammengehörige Viertelsysteme heissen. Es ist klar, dass sie folgende Eigenschaften haben:

\SourceLeaf{379}
\begin{enumerate}
\renewcommand{\labelenumi}{\arabic{enumi})}
\item Jedes Viertelsystem besteht aus $\frac14(\nu-1)=\frac14(\alpha\alpha+\beta\beta-1)$ Zahlen.
\item Das Charakteristische eines Viertelsystems ist, dass keine der darin befindlichen Zahlen weder selbst, noch ihr Product in $i,-1$, oder $-i$, einer andern aus demselben Viertelsystem congruent ist, jede durch $\mu$ nicht theilbare Zahl aber, entweder selbst oder ihr Product durch $i,-1$, oder $-i$ sich darin findet, oder einer daraus congruent ist.
\item So wie aus der Multiplication der Zahlen in $\kappa$ mit $i,-1$ und $-i$, resp.\ die Zahlen in $\kappa',\kappa'',\kappa'''$ entstehen, oder ihnen congruente, so reproducirt die Multiplication der Zahlen in $\kappa'$, mit jenen Factoren, resp.\ die Zahlen in $\kappa'',\kappa''',\kappa$; die Multiplication der Zahlen $\kappa''$ reproducirt auf ähnliche Weise die Zahlen $\kappa''',\kappa,\kappa'$; endlich die Multiplication der Zahlen $\kappa'''$ reproducirt $\kappa,\kappa',\kappa''$.
\end{enumerate}

Kürze halber kann diese gegenseitige Abhängigkeit der vier Viertelsysteme so ausgedrückt werden
\[
\begin{aligned}
\kappa'&\equiv i\kappa,\\
\kappa''&\equiv-\kappa\equiv i\kappa',\\
\kappa'''&\equiv i\kappa''\equiv-\kappa'\equiv-i\kappa.
\end{aligned}
\]
